\chapter{几何光学成像}
\problem{2-1}{以第一章第1节例题2中所用的光线追迹作图法证明图2-6中 $Q$ 和 $Q'$ 是一对共轭点。
\par\medskip\noindent
\begin{minipage}[t]{.35\linewidth}\vspace{0pt}\centering
\drawfig{1}{book_2_6}\par\smallskip
{\footnotesize 图2-6\quad 折射球面的齐明点}
\end{minipage}\hfill
\begin{minipage}[t]{.61\linewidth}\vspace{0pt}\centering
\drawfig{1}{book_2_19}\par\smallskip
{\footnotesize 图2-19\quad 用作图法求二次成像（补充参考图）}
\end{minipage}\par
\missing{题干引用的第一章第1节例题2正文；图2-6及补充图2-19已补入。}}
\answerarea
\clearpage
\problem{2-2}{证明（2.4）式。图2-6中，角度 $u=\angle MQA$、$u'=\angle MQ'A$ 与距离 $\overline{QA}$、$\overline{AQ'}$ 之间存在下列关系：
\begin{equation*}
\frac{\sin u}{\sin u'}=\frac{\overline{AQ'}}{\overline{QA}}.
\tag{2.4}
\end{equation*}
\bookfig[0.38]{book_2_6}{图2-6\quad 折射球面的齐明点}}
\answerarea
\clearpage
\problem{2-3}{根据反射定律推导球面反射镜的物像距公式（2.23）和焦距公式（2.24）。}
\answerarea
\problem{2-4}{物体放在凸球面反射镜前何处，可产生大小与物体相等的倒立实像？}
\answerarea
\clearpage
\problem{2-5}{凹面镜的半径为 $40\,\mathrm{cm}$，物体放在何处成放大两倍的实像？放在何处成放大两倍的虚像？}
\answerarea
\problem{2-6}{要把球面反射镜前 $10\,\mathrm{cm}$ 处的灯丝成像于 $3\,\mathrm m$ 处的墙上，镜形应是凸的还是凹的？半径应有多大？这时像放大了多少倍？}
\answerarea
\clearpage
\problem{2-7}{按已约定的正负号法则（I）、（II）、（III）、（IV）、（V）等，标示出本题各图中的物距 $s$、像距 $s'$、曲率半径 $r$、光线倾角 $u$、$u'$ 的绝对值。比较各图中折射率 $n$、$n'$ 的大小，指明各图中物像的虚实。\fig[0.85]{2_7}}
\answerarea
\problem{2-8}{分别根据上题各图推导球面折射成像公式（2.19）。}
\answerarea
\clearpage
\problem{2-9}{若空气中一均匀球形透明体能将平行光束会聚于其背面的顶点上，此透明体的折射率应等于多少？}
\answerarea
\problemfig[0.31]{2-10}{如本题图，一平行平面玻璃板的折射率为 $n$，厚度为 $h$，点光源 $Q$ 发出的傍轴光束（即接近于正入射的光束）经上表面反射，成像于 $Q'_1$；而折射线穿过上表面后在下表面反射，再从上表面折射的光束成像于 $Q'_2$。证明 $Q'_1$ 与 $Q'_2$ 间的距离为 $2h/n$。
\par【提示：把平面看成 $r\to\infty$ 的球面，并利用球面折射公式计算。】}{2_10}
\answerarea
\clearpage
\problemfig[0.43]{2-11}{如本题图，一会聚光束本来交于 $P$ 点，插入一折射率为 $1.50$ 的平面平行玻璃板后，像点移至 $P'$。求玻璃板的厚度。}{2_11}
\answerarea
\problem{2-12}{根据费马原理推导傍轴条件下球面反射成像公式（2.23）。}
\answerarea
\clearpage
\problem{2-13}{根据费马原理推导傍轴条件下球面折射成像公式（2.19）。}
\answerarea
\clearpage
\problem{2-14}{求联立方程式（2.15）和（2.16）的解，证明它所代表的共轭点就是第1.4节中给出的齐明点。
\begin{empheq}[left=\empheqlbrace]{align}
\frac{s^2}{n^2(s+r)^2}-\frac{s'^2}{n'^2(s'-r)^2}&=0,\tag{2.15}\\[5pt]
\frac{1}{n^2(s+r)}-\frac{1}{n'^2(s'-r)}&=0.\tag{2.16}
\end{empheq}
其中 $s,s'$ 为依原书正负号约定的物距与像距，$r$ 为球面曲率半径，$n,n'$ 分别为物方和像方介质的折射率。
\par\smallskip\textbf{齐明点：}图2-6中的 $Q,Q'$ 这对共轭点称为折射球面的齐明点（aplanatic points）。图中
\[u=\angle MQA,\qquad u'=\angle MQ'A,\qquad
\frac{\sin u}{\sin u'}=\frac{\overline{AQ'}}{\overline{QA}}\quad\text{（2.4）}.\]
\bookfig[0.38]{book_2_6}{图2-6\quad 折射球面的齐明点}}
\answerarea
\clearpage
\problem{2-15}{某透镜用 $n=1.50$ 的玻璃制成，它在空气中的焦距为 $10.0\,\mathrm{cm}$，它在水中的焦距为多少？（水的折射率为 $4/3$。）}
\answerarea
\problem{2-16}{一薄透镜折射率为 $1.50$，光焦度为 $5.00\,\mathrm D$。将此镜浸入某液体，光焦度变为 $-1.00\,\mathrm D$，求此液体的折射率。}
\answerarea
\clearpage
\problemfig[0.27]{2-17}{用一曲率半径为 $20\,\mathrm{cm}$ 的球面玻璃和一平面玻璃黏合成空气透镜，将其浸于水中（见本题图）。设玻璃壁厚可忽略，水和空气的折射率分别为 $4/3$ 和1，求此透镜的焦距。此透镜是会聚的还是发散的？}{2_17}
\answerarea
\problem{2-18}{一凸透镜的焦距为 $12\,\mathrm{cm}$，填充下表中的空白，并作出相应的光路图。\lenstable}
\answerarea
\clearpage
\problem{2-19}{一凹透镜的焦距为 $12\,\mathrm{cm}$，填充下表中的空白，并作出相应的光路图。\lenstable}
\answerarea
\problem{2-20}{在 $5\,\mathrm{cm}$ 焦距的凸透镜前放一小物，要虚像成像于 $25\,\mathrm{cm}$ 到无穷远之间，物应放在什么范围？}
\answerarea
\clearpage
\problem{2-21}{一光源与屏间的距离为 $1.6\,\mathrm m$，用焦距为 $30\,\mathrm{cm}$ 的凸透镜插在二者之间，透镜应放在什么位置，才能使光源成像于屏上？}
\answerarea
\problem{2-22}{屏幕放在距物 $100\,\mathrm{cm}$ 远处，二者之间放一凸透镜。当前后移动透镜时，我们发现透镜有两个位置可以使物成像在屏幕上。测得这两个位置之间的距离为 $20.0\,\mathrm{cm}$，求：（1）这两个位置到幕的距离和透镜的焦距；（2）两个像的横向放大率。}
\answerarea
\clearpage
\problem{2-23}{如上题，在固定的物与幕之间移动凸透镜。证明：要使透镜有两个成像位置，物和幕之间的距离必须大于四倍焦距。}
\answerarea
\problemfig[0.41]{2-24}{如本题图，$L_1$、$L_2$ 分别为凸透镜和凹透镜。前放一小物，移动屏幕到 $L_2$ 后 $20\,\mathrm{cm}$ 的 $S_1$ 处接到像。现将凹透镜 $L_2$ 撤去，将屏移前 $5\,\mathrm{cm}$ 至 $S_2$ 处，重新接收到像。求凹透镜 $L_2$ 的焦距。}{2_24}
\answerarea
\clearpage
\problem{2-25}{一光学系统由一焦距为 $5.0\,\mathrm{cm}$ 的会聚透镜 $L_1$ 和一焦距为 $10.0\,\mathrm{cm}$ 的发散透镜 $L_2$ 组成，$L_2$ 在 $L_1$ 之右 $5.0\,\mathrm{cm}$。在 $L_1$ 之左 $10.0\,\mathrm{cm}$ 处放置一小物，求经此光学系统后所成的像的位置和横向放大率。用作图法验证计算结果。}
\answerarea
\problem{2-26}{用逐次成像法解第4.4节中的例题4，并将结果与之比较。
\bookfig[0.94]{book_2_29}{图2-29\quad 第4.4节例题4配图}
\missing{例题4的完整题干与已知参数；图2-29已补入。}}
\answerarea
\clearpage
\problem{2-27}{用作图法求本题各图中 $PQ$ 的像（入射线从左到右）。\fig[0.9]{2_27}}
\answerarea
\problem{2-28}{用作图法求本题图中 $PQ$ 的像（入射线从左到右）。\fig[0.7]{2_28}}
\answerarea
\clearpage
\problem{2-29}{用作图法求本题图中联合光具组的主面和焦点。\fig[0.75]{2_29}}
\answerarea
\problem{2-30}{用作图法求本题图中正方形 $ABCD$ 的像。\fig[0.65]{2_30}}
\answerarea
\clearpage
\problem{2-31}{验算第7.7节图2-63a和b中给出的两种拉姆斯登目镜的主面与焦点。
\par\smallskip 向场镜 $L_1$ 和接目镜 $L_2$ 的焦距和间隔之比，在图（a）中为
\[f'_1:f'_2:d=1:1:1.\]
图中 $F,F'$ 为整个系统的焦点，$H,H'$ 为主点，$\mathcal H,\mathcal H'$ 表示相应的主平面；物体 $QP$ 放在焦点 $F$ 内侧附近。图（b）为缩短两镜间距的结构示意。
\bookfig[0.98]{book_2_63}{图2-63\quad 拉姆斯登目镜}}
\answerarea
\problem{2-32}{求下表中厚透镜的焦距和主面、焦点的位置，并作图表示。已知玻璃的折射率为 $1.500$，两界面顶点间的距离为 $1.00\,\mathrm{cm}$，透镜放在空气中。
\begin{center}\begin{tabular}{c|c|r|r}\hline &形状&$r_1/\mathrm{cm}$&$r_2/\mathrm{cm}$\\\hline
（1）&双凸&10.0&$-10.0$\\（2）&凸凹&10.0&20.0\\（3）&凹凸&$-15.0$&$-20.0$\\\hline\end{tabular}\end{center}}
\answerarea
\clearpage
\problem{2-33}{求放在空气中玻璃球的焦距和主面、焦点的位置，并作图表示。已知玻璃球的半径为 $2.00\,\mathrm{cm}$，折射率为 $1.500$。}
\answerarea
\problem{2-34}{上题中玻璃球表面上有一斑点，计算从另一侧观察此斑点像的位置和放大率，并用作图法验证之。}
\answerarea
\clearpage
\problemfig[0.42]{2-35}{（1）用作图法求本题图中光线1的共轭线；（2）在图上标出光具组节点 $N$、$N'$ 的位置。}{2_35}
\answerarea
\problemfig[0.42]{2-36}{已知 $1-1'$ 是一对共轭光线，求光线2的共轭线（见本题图）。}{2_36}
\answerarea
\clearpage
\problem{2-37}{对一光具组，测得当物距改变 $\Delta x$ 时，像距改变 $\Delta x'$，同时横向放大率由 $V_1$ 变为 $V_2$。试证明此光具组的焦距为
\[f=\frac{\Delta x}{\dfrac1{V_2}-\dfrac1{V_1}},\qquad f'=\frac{\Delta x'}{V_2-V_1}.\]
这里提供了一种测焦距的方法，它与测焦点位置的方法配合起来，可以确定光具组的主面。}
\answerarea
\problem{2-38}{一架幻灯机的投影镜头的焦距为 $7.5\,\mathrm{cm}$，当幕由 $8\,\mathrm m$ 移至 $10\,\mathrm m$ 远时，镜头移动多少距离？}
\answerarea
\clearpage
\problem{2-39}{某照相机与拍摄物体的最近距离为 $1\,\mathrm m$，装上光焦度为 $2\,\mathrm D$ 的近视镜片，能拍摄的最近距离为多少？（假设近拍镜与照相机镜头是密接的。）}
\answerarea
\problem{2-40}{某人对 $2.5\,\mathrm m$ 以外的物看不清，需配多少度的眼镜？另一人对 $1\,\mathrm m$ 以内的物看不清，需配怎样的眼镜？}
\answerarea
\clearpage
\problem{2-41}{计算 $2\times$、$3\times$、$5\times$、$10\times$ 放大镜或目镜的焦距。}
\answerarea
\problem{2-42}{一架显微镜，物镜焦距为 $4\,\mathrm{mm}$，中间像放在物镜像方焦点后面 $160\,\mathrm{mm}$ 处，如果目镜是 $20\times$ 的，显微镜的总放大率是多少？}
\answerarea
\clearpage
\problem{2-43}{一架显微镜的物镜和目镜相距 $20.0\,\mathrm{cm}$，物镜焦距为 $7.0\,\mathrm{mm}$，目镜焦距为 $5.0\,\mathrm{mm}$，把物镜和目镜都看成是薄透镜。（1）求被观测物到物镜的距离；（2）物镜的横向放大率；（3）显微镜的总放大率；（4）焦深。}
\answerarea
\problem{2-44}{物镜、目镜皆为会聚的望远镜称为开普勒型望远镜（图2-29），物镜会聚而目镜发散的望远镜，称为伽里略型望远镜。
\begin{enumerate}\item 画出伽里略型望远镜的光路。\item 一伽里略型望远镜的物镜和目镜相距 $12\,\mathrm{cm}$，若望远镜的放大率为 $4\times$，物镜和目镜的焦距各多少？\end{enumerate}}
\answerarea
\clearpage
\problem{2-45}{拟制一个 $3\times$ 的望远镜，已有一个焦距为 $50\,\mathrm{cm}$ 的物镜，问，在（1）开普勒型和（2）伽里略型中，目镜的光焦度以及物镜、目镜的距离各多少？}
\answerarea
\problem{2-46}{一望远镜的物镜直径为 $5.0\,\mathrm{cm}$，焦距为 $20\,\mathrm{cm}$，目镜直径为 $1.0\,\mathrm{cm}$，焦距为 $2.0\,\mathrm{cm}$，求此望远镜的入射光瞳和出射光瞳的位置和大小。}
\answerarea
\clearpage
\problem{2-47}{望远镜的孔径光阑和入射光瞳通常就是其物镜的边缘。求出射光瞳的位置，并证明出射光瞳直径 $D'$ 与物镜直径 $D_0$ 之比为
\[D'=\frac{D_0}{|M|},\]
其中 $M=-f'_{\mathrm{ob}}/f'_{\mathrm{oc}}$ 是望远镜的视角放大率。}
\answerarea
\problem{2-48}{将望远镜倒过来可作激光扩束之用。设一望远镜物镜焦距为 $30\,\mathrm{cm}$，目镜焦距为 $1.5\,\mathrm{cm}$，它能使激光光束的直径扩大几倍？}
\answerarea
\clearpage
\problem{2-49}{显微镜的孔径光阑和入射光瞳通常就是其物镜的边缘。求：（1）出射光瞳的位置；（2）证明在傍轴近似下出射光瞳的直径 $D'$ 与入射孔径角 $u_0$ 的关系是
\[D'=\frac{2s_0n u_0}{|M|},\]
式中 $s_0=25\,\mathrm{cm}$ 是明视距离，$M$ 是显微镜的视角放大率，$n$ 是物方折射率（除油浸物镜外，$n\approx1$）。严格的公式应把 $u_0$ 换成 $\sin u_0$，参见习题2-56。\missing{原截图脚注的“2-56”疑指本批2-53，保留原引用待核。}}
\answerarea
\problem{2-50}{如本题图中 $L_1$、$L_2$ 是两个会聚透镜，$Q$ 是物点，$DD$ 是光阑。已知焦距 $f_1=2a$，$f_2=a$，图中标示各距离为 $s=10a$、$l=4a$、$d=6a$；此外透镜与光阑半径之比是 $r_1=r_2=3r_3$。求此光具组的孔径光阑、入射光瞳、出射光瞳、入射窗和视场光阑的位置和大小。\fig[0.9]{2_50}}
\answerarea
\clearpage
\problemfig[0.39]{2-51}{惠更斯目镜的结构详见第7.7节图2-62。如本题图所示，今在两透镜间放一光阑 $AA$，设透镜 $L_1$、$L_2$ 和光阑的直径分别为 $D_1$、$D_2$ 和 $D$。试证：
\begin{enumerate}\item 向场镜 $L_1$ 成为孔径光阑的条件为 $D_1<D_2$；\item 光阑 $AA$ 成为视场光阑的条件是 $D/2<D_2$；\item 这时对 $F'_2$ 点计算的出射孔径角 $u'_0$ 由下式确定：$\tan u'_0=D_1/(2a)$。\end{enumerate}}{2_51}
\answerarea
\problem{2-52}{证明在一对齐明点附近不可能有另一对齐明点。}
\answerarea
\clearpage
\problem{2-53}{试证明，显微镜出射光瞳直径 $D'$ 由下式确定：
\[D'=\frac{2s_0n\sin u_0}{|M|}=2s_0\frac{\mathrm{N.A.}}{|M|},\]
即 $D'$ 正比于数值孔径 $\mathrm{N.A.}=n\sin u_0$，反比于视角放大率 $M$。（$s_0=25\,\mathrm{cm}$ 是明视距离。）}
\answerarea
\problem{2-54}{拟用冕牌玻璃 K9（$n_D=1.5163$，$n_F=1.5220$，$n_C=1.5139$）和重火石玻璃 F4（$n_D=1.6199$，$n_F=1.6321$，$n_C=1.6150$）来做消色差胶合透镜，焦距为 $100\,\mathrm{mm}$。若已确定其负透镜的非黏合面为平面，试求其余各面的曲率半径。}
\answerarea
\clearpage
\problem{2-55}{太阳表面的辐射亮度为 $2\times10^7\,\mathrm{W/(m^2\cdot sr)}$，用相对孔径 $D/f=1.5$ 的放大镜将阳光聚焦成光斑的最大辐射照度为多少？}
\answerarea
\problem{2-56}{一屏放在离烛 $100\,\mathrm{cm}$ 处，把一会聚薄透镜放在烛和屏之间，透镜有两个位置可以在屏上得到烛的像，这两个位置相距 $20\,\mathrm{cm}$，在屏上烛的两个像的照度相差多少倍？}
\answerarea
\clearpage
\problem{2-57}{望远镜物镜的直径为 $75\,\mathrm{mm}$，当放大率为（1）20倍，（2）25倍，（3）50倍时，求望远镜中月亮的像的主观亮度与天然主观亮度之比。眼瞳瞳孔的直径为 $3.0\,\mathrm{mm}$。}
\answerarea
\problem{2-58}{一天文望远镜的物镜直径等于 $18\,\mathrm{cm}$，透光系数为 $0.50$。已知肉眼可直接观察到六等星。求：
\begin{enumerate}\item 用此望远镜所能看到的最弱星等；\item 最适宜观察星的放大率（正常放大率）；\item 当放大率为10倍时可以见到星的等次。\end{enumerate}
设眼睛瞳孔的直径可取 $3.0\,\mathrm{mm}$。【注：星等增加一等，其亮度减小 $\sqrt[5]{100}\approx2.5$ 倍。】}
\answerarea
\clearpage
\problem{2-59}{求数值孔径为 $1.5$ 的显微镜的正常放大率，设瞳孔直径为 $3.0\,\mathrm{mm}$。}
\answerarea
\clearpage
